\subsection*{A physical copy distribution and a learned whole-span action}
Let $x_i$ denote the exact bitmap at source position $i$, $g_t$ a sigmoid writing probability, $a_{ti}$ normalized attention over source positions and $b_t(x)$ the distribution over the 7,990 printed bitmaps and stop. We used
\[
  p_t(x)=g_t b_t(x)+(1-g_t)\sum_{i:x_i=x}a_{ti}.
\]
The sum merges physically identical outcomes before greedy decoding. The question precedes the passage in the causal context. A learned local convolution supplies contextual keys for the source positions. An adjacency term favours the next source position when the previous emitted bitmap equals its predecessor. The selected bitmap, including an out-of-bank copied form, supplies the next visual input. A stop decision can only come from the printed branch. When the training answer occurs as a complete span, an auxiliary first-step loss of weight 0.35 rewards its possible starts. This design adapts the copy mechanism of \href{https://aclanthology.org/P17-1099/}{See, Liu and Manning (2017)} to exact bitmap identity.

We fitted the top four causal blocks, last normalization and the new 246,914-parameter pointer on 2,565 DuReader records, 1,918 paired CMRC questions and 2,048 captions on 512 training plates. Each of three arms trained on two passes, 818 updates and 248,356 answer positions. The genuine fit took 67.9 seconds and 3.23 GiB peak allocated memory on one RTX 4090 D, excluding previous language pretraining and interface calibration. The two controls held target batches fixed, with uniform source positions or questions exchanged on each paired training page. Both natural tests used 128 passages with two distinct human answers. DRCD used eligible records from the official test file. Upstream CKIP benchmark exposure was not audited. The 48 historical evaluation cards used held-out families, two headwords per card and two stage orders per headword.

A frozen CKIP-based visual reader scores all nonempty source spans up to 48 cells by the sum of their endpoint logits. A 148,033-parameter binary classifier reads its summary vector, mean selected-span state, span score and margin, and length features. The target is physical containment of the complete training answer. A score of at least one half selects the span as a single terminal action. Otherwise the frozen causal writer runs to its own stop or limit. The original selector used eight passes and a shuffled-label comparison. Follow-up fits added either answer-containing sentences or random excerpts of identical lengths from the same training passages. The random windows contained 174 answers. The sentence windows contained all 1,738. The follow-up compares joint changes in support and sampling and cannot isolate the effect of length alone. The 16 additional short scenes, with exchanged names, were authored for this evaluation and excluded from all fits. For retrieval, a separate ranker used 591 training questions with both supporting and nonsupporting candidate windows among 1,536 sampled questions. Image generation, span return, target comparison and all evaluation figures used the physical glyph sequences.
